% abstract.tex : Hebrew abstract for ACO-SLAM paper
% Compiler: XeLaTeX
% Language: Hebrew (primary)

\selectlanguage{hebrew}

מאמר זה מציג את 
\en{ACO-SLAM},
 מסגרת חדשנית לניווט אוטונומי של להקות רחפנים למשימות חיפוש והצלה (
\en{SAR}
) בסביבות ללא 
\en{GPS}.
 המסגרת משלבת אופטימיזציית מושבת נמלים (
\en{ACO}
) עם 
\en{SLAM}
 הסתברותי מבוסס גריד תפוסה, תוך שימוש באופטימיזציית גרף תנוחות מבוזרת עם קונצנזוס בהשראת נמלים. מיזוג חיישנים רב-מודאלי מבוסס פרומונים משלב 
\en{LiDAR},
 מצלמה ו-
\en{IMU}
 תוך שקלול דינמי של מהימנות כל חיישן. תכנון מסלול אדפטיבי מאזן בין חקר מרחבי לבין ניצול מידע קיים באמצעות קצב התאיידות פרומונים משתנה בזמן. תוצאות סימולציה מקיפות בסביבת 
\en{Gazebo/Unity}
 עם תרחיש 
\en{SAR}
 עירוני בגודל 100×100 מטר ו-5 רחפנים מראות שיפור של 34\% בשגיאת מסלול מוחלטת (
\en{ATE}
) לעומת 
\en{Swarm-SLAM}
 (5.4 ס"מ לעומת 8.2 ס"מ), שיפור של 8.5\% בכיסוי המרחבי (94.8\% לעומת 87.3\%), וקיצור של 23.7\% בזמן השלמת המשימה (23.2 דקות לעומת 30.4 דקות). ניסויי מעבדה עם 3 רחפנים בסביבה פנימית מבוקרת אישרו את מגמות הסימולציה. המסגרת המוצעת מהווה צעד משמעותי לקראת פריסה מבצעית של להקות רחפנים אוטונומיות למשימות חיפוש והצלה בסביבות מורכבות ומשתנות.